\documentclass[12pt, ukrainian]{article}
\usepackage[a4paper]{geometry}
\setlength\parindent{0mm}
\geometry{top=1.5cm,bottom=1.5cm,left=1.3cm,right=1.5cm}
\pagestyle{empty}
\usepackage[T2A]{fontenc}
\usepackage[utf8]{inputenc}
\usepackage[ukrainian]{babel}
\usepackage{graphicx}
\usepackage{caption}
\usepackage{caption}
\DeclareCaptionFormat{nocaption}{}
\captionsetup{format=nocaption,aboveskip=0pt,belowskip=0pt}
\usepackage[Export]{adjustbox}
\adjustboxset{max size={0.9\linewidth}{0.9\paperheight}}
\usepackage{verbatim, listings, clrscode3e, fancyvrb}
\def\code#1{\Verb!#1!}
\begin{document}
%begin
\begin {large}

\textbf{22.11.2024 Контрольна робота, варіант  \No { 54 }}\par


{

\begin{center}
\textbf{Завдання 1}
\end{center}

По яких днях тижня відбувались лекції з предмету?


\begin{center}
\textbf{Завдання 2}
\end{center}

Як зовуть (ім'я, по батькові) викладача, що вів практичні заняття?


\begin{center}
\textbf{Завдання 3}
\end{center}

Нехай int var=13; Один потік збільшив змінну var на 3, а інший збільшив її в три рази.

гонка даних (data race), стан гонки (race condition)

Виберіть усі правильні твердження (якщо такі є).\par
\vspace{2mm}
(\,1\,) Тут гарантовано нема гонки даних, але може бути стан гонки.%bad

(\,2\,) Тут може бути тільки стан гонки.%bad

(\,3\,) Тут може не бути стану гонки.%bad

(\,4\,) Тут гарантовано є стан гонки.%good


\begin{center}
\textbf{Завдання 4}
\end{center}

Виберіть усі правильні твердження (якщо такі є).\par
\vspace{2mm}

(\,1\,) Коли потік спить, він не споживає ресурс процесора.%good

(\,2\,) Довільний потік може розблокувати заблокований м’ютекс.%bad

(\,3\,) Виклик \code{this\_thread::sleep\_for(5s)} змушує потік проспати рівно 5 секунд.%bad

(\,4\,) Об’єкти класів promise та future можна використовувати тільки в багатопотокових програмах.%bad


\begin{center}
\textbf{Завдання 5}
\end{center}

Виберіть усі правильні твердження (якщо такі є).\par
\vspace{2mm}

(\,1\,) Мова С++20 має стандартну бібліотеку вільних від блокувань структур даних.%bad

(\,2\,) Операція над вільною від блокувань структурою даних може виконуватись скільки завгодно довго.%good

(\,3\,) Операція над вільною від блокувань структурою даних завжди виконається за обмежений час.%bad

(\,4\,) Вільна від блокувань структура даних гарантує, що жоден потік не буде очікувати доступу до виконання операцій над цією структурою даних.%bad
\newpage


\begin{center}
\textbf{Завдання 6}
\end{center}

Виберіть усі правильні твердження (якщо такі є).\par
\vspace{2mm}

(\,1\,) Потік може мати доступ до автоматичних змінних іншого потоку і це безпечно.%bad

(\,2\,) Динамічні змінні не можуть розділятися між потоками.%bad

(\,3\,) Потік має доступ тільки до своїх власних даних та до глобальних статичних даних.%bad

(\,4\,) Потік може мати доступ до автоматичних змінних іншого потоку, але це небезпечно.%good


\begin{center}
\textbf{Завдання 7}
\end{center}

У цьому фрагменті коду
\begin{verbatim}
template <class T> void f(T a){}

void g(){
    int i=5;
    f(unique_ptr<int>(new int(i)))+(unique_ptr<int>(new int(35)));
\end{verbatim}
(\,1\,) є невизначена поведінка%bad

(\,2\,) є стан гонки%bad

(\,3\,) є гонка даних%bad

(\,4\,) можливий виток ресурсів%good

}
\end {large}

\vspace{4mm}
%end
\newpage
\end{document}

